\documentclass[10pt,a4paper]{amsart}
\usepackage[margin=19mm]{geometry}
\usepackage{amsmath,amssymb,mathtools}
\usepackage[scr=rsfs,cal=euler]{mathalfa}
\usepackage{xfrac}
\usepackage[unicode,hidelinks,bookmarks=false]{hyperref}
\newcommand{\ie}{i.e.\ }
\newcommand{\cf}{cf.\ }
\newcommand{\resp}{respectively\ }
\newcommand{\Subsection}{\subsection}
\providecommand{\nobreakdash}{\nobreak\hbox{-}}
\setlength{\emergencystretch}{2em}
\multlinegap0pt
\usepackage{amssymb}
\usepackage{tikz}
\usepackage{mathtools}
\newtheorem{thm}{Theorem}
\newtheorem{prop}{Proposition}
\newcommand{\GL}{\mathrm{GL}}
\newcommand{\Irr}{\operatorname{Irr}}
\newcommand{\Sp}{\operatorname{Sp}}
\newcommand{\cT}{\mathcal{T}}
\newcommand{\cc}{\mathrm{c}}
\def\rO{{\mathrm{O}}}
\newcommand{\rU}{\mathrm{U}}
\title{The theta correspondence over finite fields}
\author{Anne-Marie Aubert}
\date{Source: arXiv:2603.25658v1 (CC BY 4.0)}
\begin{document}
\maketitle

\noindent\emph{Source and licence.} The theta correspondence over finite fields. Version arXiv:2603.25658v1 (CC BY 4.0), distributed by arXiv under the Creative Commons Attribution 4.0 International licence, http://creativecommons.org/licenses/by/4.0/, as recorded in the arXiv licence metadata for this exact version and shown on the abstract page https://arxiv.org/abs/2603.25658v1. Full licence notice: \texttt{licenses/arxiv-2603.25658-CC-BY-4.0.txt}.\par\smallskip

\noindent\textit{Editorial scope.} Counted unit: the article's prop pro:principal-unipotent (source0.tex lines 1400--1404). The article's complete original proof of it is delivered with it (source0.tex lines 1405--1407, one \texttt{proof} environment of 3 lines). No mathematics is written, repaired or re-derived: the statement and the proof are the article's own lines, in the article's own notation, and no retained line is substituted or re-flowed.  Retained uncounted as the source-local context the proof needs: lines 879--890; lines 1349--1364.  Nothing was omitted from any retained span, so the package carries no omission record.  No retained line contains a citation, so the wrapper carries no bibliography and no bibliography entry.  No retained line contains a figure environment, an image-inclusion command or drawing code, so no image is delivered and the package declares no asset dependency.  Editorial formatting only, in addition: an AMS \texttt{amsart} preamble that restates the article's own declarations and macros used by the retained lines, namely the environments \texttt{thm}, \texttt{prop} on separate counters, together with the standard packages those lines need.  The retained environments print in this document as Theorem 1, Proposition 1 in order of appearance, whereas the article prints the same items with the numbering of its own full document.  The complete native source of the article as distributed in the arXiv e-print for this exact version is bundled unchanged as \texttt{sources/197233/source0.tex}.
\begin{thm} \label{thm:HC-HC}
Let $(G,G')=(G_n,G'_{n'})$ be a reductive dual pair, let $\sigma\in\Irr_\cc(L)$, and let $\pi\in\Irr^\sigma(G)$. 

Then every $\pi'\in\theta_{G'}(\pi)$ belongs to $\Irr^{\sigma'}(G')$, where
\begin{enumerate} 
\item if $\overline{n}\le\overline{n}'$
\[\sigma'=\theta^0(\pi_{n_0})\otimes\tau_1\otimes\cdots\otimes\tau_r\otimes 1\otimes\cdots\otimes 1;\]
\item otherwise there exists a sequence $i_1$, $\ldots$, $i_t$ such that
\[\sigma'=\theta^0(\pi_{n_0})\otimes\tau_1\otimes\cdots\otimes\widehat{\tau_{i_1}}\otimes\cdots\otimes\widehat{\tau_{i_t}}\otimes\cdots\otimes\tau_r,\]
\end{enumerate}
with $\theta^0(\pi_{n_0})$ the first  occurrence of $\pi_{n_0}$ with respect to the Witt tower associated to $G'$.
\end{thm}

\begin{thm} \label{thm:ucuspidals} \

\begin{enumerate}
\item We suppose that  $G=\Sp_{2n}(k)$ and $G'=\rO_{2n'}^\epsilon(k)$. The following holds:
\begin{enumerate}
\item If $\pi$ is a cuspidal irreducible  representation of $G$, then $\theta_{n'_{\cT'}(\pi)}(\pi)$ is a singleton $\{\pi'\}$ with $\pi'$ cuspidal irreducible.
\item If $\pi\in\Irr(G_n)$ is unipotent then any $\pi'\in \theta_{G'}(\pi)$ is unipotent.
\item The representation $\pi_c^G$ of $G=\Sp_{2(c^2+c)}(k)$ corresponds to the representation $\pi^{G'}_{c',-}$ with $c'=c$ if $\epsilon$ is the sign of $(-1)^c$, and to the representation $\pi^{G'}_{c',+}$ with $c'=c+1$ otherwise.
\end{enumerate}
\item We suppose that  $G=\rU_{n}^\epsilon(k)$ and $G'=\rU_{n'}^{\epsilon'}(k)$. The representation $\pi_c^G$ of $G$ corresponds to the representation $\pi^{G'}_{c'}$, where
\begin{enumerate}
\item if $c=0$, then $c'=0$;
\item if $c\ne 0$, then $c'=c\pm 1$ such that $\epsilon'=(-1)^{\frac{c'(c'+1)}{2}}$.
\end{enumerate}
\end{enumerate}
\end{thm}

\begin{prop} \label{pro:principal-unipotent}
For the  irreducible reductive dual pairs 
\[(G,G')\in\{(\Sp_{2n}(k),\rO_{2n'}^+(k)),(\rU_{n}^+(k),\rU_{n}^+(k)),(\GL_n(k),\GL_{n'}(k))\},\] 
and only for these ones,  the images by the theta correspondence of every unipotent representation in the principal series of $G$ belongs to the (unipotent) principal series of $G'$.
\end{prop}

\begin{proof}
It follows from the combination of Theorem~\ref{thm:HC-HC} with Theorem~\ref{thm:ucuspidals}.
\end{proof}

\end{document}
